Small paired evaluations report conclusions, and a conclusion is a statement
about how a handful of answers happened to be graded. This paper takes one
completed question-answering evaluation --- \Questions{} questions answered
under \ConditionsWord{} retrieval conditions, \Answers{} graded answers in total --- and
asks of each of its \Contrasts{} recorded contrasts the same arithmetic
question: what is the fewest individual gradings that would have to have come
out the other way for the reported conclusion to stop holding? Because a paired
contrast reduces to four counts and one grading moves a question between two of
them, the answer is found by exhaustive search and is a minimum rather than a
bound. Three conclusions were named in advance. The finding that the multimodal
condition trails the text-only one on image-dependent questions rests on
\HarmK{} grading: \HarmCandidates{} of the \HarmGradings{} individual gradings
on that subset would each on their own erase it. That flip is not hypothetical.
An independent regrade of the same answers by \Graders{} graders finds exactly
\RegradeContested{} grading they disagree about, and applying it as the record
states moves the difference from \HarmDiff{} to \RegradedDiff{} percentage
points. The second conclusion, a \HelpDiff{}-point difference in the other
direction, needs \HelpK{} gradings to reverse --- but its non-significant
verdict needs \HelpVerdictK{}, so the direction of that result is \HelpRatio{}
times better supported than the verdict attached to it. The third, the only
contrast to survive correction for the \Contrasts{} comparisons, needs \SigK{}.
Two further facts about the same rows sharpen the point: \Incapable{} of the
\Contrasts{} contrasts had a smallest attainable p-value above the threshold and
so could not have been called significant however the answers came out, and
\AtFloor{} returned exactly the most extreme value available to them. None of this
adjudicates whether the retrieval conditions differ, and no claim that they do
or do not is made here. What it establishes is narrower and checkable: at this
scale a conclusion is a property of individual gradings, the direction of a
result and the verdict on it are separately fragile and can differ by a factor
of \HelpRatio{}, and reporting either without the count of gradings it rests on
presents an arithmetic accident as a finding.
